% !TeX root = ../main.tex
% ============================================================
% LÁMINA 5 — ESTÁNDARES CLÍNICOS (fondo claro)
% ============================================================
\begin{frame}[plain]
  \begin{tikzpicture}[remember picture, overlay]

    \fill[paperBg] (current page.south west) rectangle (current page.north east);

    \node[anchor=north west, align=left]
      at ([xshift=1.4cm, yshift=-0.55cm]current page.north west)
      {\kicker[accentBlue]{04 \,$\cdot$\, Estándares clínicos}};

    \node[anchor=north west, align=left, text width=13.2cm]
      at ([xshift=1.4cm, yshift=-1.05cm]current page.north west)
      {\Large\bfseries\sffamily\color{inkText}Terminologías e interoperabilidad como base del sistema};

    % --- 3 tarjetas oscuras en fila ---
    \def\stdW{4.2cm}
    \def\stdH{4.5cm}
    \def\stdGap{0.3cm}
    \def\stdTextW{3.5cm}
    \def\stdY{-2.1cm}

    % Tarjeta 1 — SNOMED CT
    % Cada bloque de texto se ancla debajo del anterior (altura variable),
    % así una línea extra por wrap nunca pisa el bloque siguiente.
    \node[anchor=north west, fill=inkDark, rounded corners=6pt,
          minimum width=\stdW, minimum height=\stdH, inner sep=0pt]
      (s1) at ([xshift=1.4cm, yshift=\stdY]current page.north west) {};
    \fill[accentTeal, rounded corners=2pt]
      ([xshift=0cm]s1.north west) rectangle ([yshift=-0.12cm]s1.north east);
    \node[anchor=north west, text width=\stdTextW, align=left] (s1title)
      at ([xshift=0.4cm, yshift=-0.55cm]s1.north west)
      {\normalsize\bfseries\sffamily\color{accentTeal}SNOMED CT};
    \node[anchor=north west] (s1sub)
      at ([yshift=-0.12cm]s1title.south west)
      {\parbox[t]{\stdTextW}{\raggedright\fontsize{8.5}{10}\selectfont\sffamily\color{white}%
       Terminología clínica de referencia}};
    \node[anchor=north west]
      at ([yshift=-0.18cm]s1sub.south west)
      {\parbox[t]{\stdTextW}{\raggedright\fontsize{8}{9.5}\selectfont\sffamily\color{mutedText}%
       Identifica y codifica con precisión semántica los conceptos médicos extraídos del dictado}};

    % Tarjeta 2 — CIE-10-ES
    \node[anchor=north west, fill=inkDark, rounded corners=6pt,
          minimum width=\stdW, minimum height=\stdH, inner sep=0pt]
      (s2) at ([xshift=1.4cm+\stdW+\stdGap, yshift=\stdY]current page.north west) {};
    \fill[accentTeal, rounded corners=2pt]
      ([xshift=0cm]s2.north west) rectangle ([yshift=-0.12cm]s2.north east);
    \node[anchor=north west, text width=\stdTextW, align=left] (s2title)
      at ([xshift=0.4cm, yshift=-0.55cm]s2.north west)
      {\normalsize\bfseries\sffamily\color{accentTeal}CIE-10-ES};
    \node[anchor=north west] (s2sub)
      at ([yshift=-0.12cm]s2title.south west)
      {\parbox[t]{\stdTextW}{\raggedright\fontsize{8.5}{10}\selectfont\sffamily\color{white}%
       Clasificación diagnóstica oficial}};
    \node[anchor=north west]
      at ([yshift=-0.18cm]s2sub.south west)
      {\parbox[t]{\stdTextW}{\raggedright\fontsize{8}{9.5}\selectfont\sffamily\color{mutedText}%
       Traduce los conceptos \textit{SNOMED CT} a códigos que se usan en la práctica clínicas españolas}};

    % Tarjeta 3 — HL7 FHIR
    \node[anchor=north west, fill=inkDark, rounded corners=6pt,
          minimum width=\stdW, minimum height=\stdH, inner sep=0pt]
      (s3) at ([xshift=1.4cm+2*\stdW+2*\stdGap, yshift=\stdY]current page.north west) {};
    \fill[accentTeal, rounded corners=2pt]
      ([xshift=0cm]s3.north west) rectangle ([yshift=-0.12cm]s3.north east);
    \node[anchor=north west, text width=\stdTextW, align=left] (s3title)
      at ([xshift=0.4cm, yshift=-0.55cm]s3.north west)
      {\normalsize\bfseries\sffamily\color{accentTeal}HL7 FHIR};
    \node[anchor=north west] (s3sub)
      at ([yshift=-0.12cm]s3title.south west)
      {\parbox[t]{\stdTextW}{\raggedright\fontsize{8.5}{10}\selectfont\sffamily\color{white}%
       Estándar de interoperabilidad}};
    \node[anchor=north west]
      at ([yshift=-0.18cm]s3sub.south west)
      {\parbox[t]{\stdTextW}{\raggedright\fontsize{8}{9.5}\selectfont\sffamily\color{mutedText}%
       Estructura el informe final para su intercambio con sistemas de Historia Clínica Electrónica}};

    % --- Nota inferior ---
    \node[anchor=north west]
      at ([xshift=1.4cm, yshift=-7.0cm]current page.north west)
      {\parbox[t]{13.2cm}{\raggedright\footnotesize\sffamily\color{mutedText}%
       Los tres estándares actúan de forma encadenada: \textit{SNOMED CT} identifica el concepto,
       \textit{CIE-10-ES} lo clasifica y \textit{HL7 FHIR} lo transporta.}};

  \end{tikzpicture}
  \end{frame}
